\section{Background and Related Work}
\label{sec:related}

\paragraph{LLM-based program repair.} Modern repair pipelines prompt an LLM
with buggy code and, usually, an observed failure (test output, stack trace),
and validate candidate patches against the test suite. Benchmarks span
algorithmic bugs (QuixBugs~\cite{quixbugs}), function-level edits
(HumanEvalFix~\cite{octopack}), synthetic mutations
(MBPP-derived~\cite{mbpp}), and real repository bugs
(BugsInPy~\cite{bugsinpy}, PyBugHive~\cite{pybughive}). Our study uses all
four regimes deliberately: the first three give executable, well-powered
comparisons; the last tests external validity.

\paragraph{Retrieval-augmented repair.} Several systems retrieve historical
bug-fix examples to condition repair. InferFix~\cite{inferfix} retrieves by
static-analyzer bug type; ReCode~\cite{recode} retrieves by algorithm type;
RAP-Gen~\cite{rapgen} learns a hybrid retriever over CodeT5. Our retrieval
signal differs: we extract a structured state from the \emph{dynamic} test
failure (runtime signal), which is orthogonal to static bug types, and we make
the reranking function explicit and ablatable rather than learned
end-to-end. Direct numeric comparison with these systems is infeasible
(closed source, different languages); our contribution is not a better
system but a controlled analysis of whether the retrieval component can pay
off at all.

RAGFix~\cite{ragfix} retrieves Stack Overflow posts through a
question-answering knowledge graph and reports repair gains on HumanEvalFix
with Llama-3 models --- a different knowledge source (crowd discussion rather
than bug-fix pairs), but the same underlying bet that retrieved context
improves repair. In a recomputation from its released per-example outputs,
the reported 70B improvement is consistent with a post-processing
(import-fixing) artifact rather than retrieval itself; a controlled ablation
on identical outputs would be needed to confirm this, so we flag it as a
caveat rather than a claim. Our own de-confounded measurements
(Section~\ref{sec:results-conversion}) agree with the direction that
retrieval is neutral-to-harmful for capable models on this benchmark.

\paragraph{Negative and analytical results on RAG.} Outside repair, RAG
studies increasingly report conditions under which retrieval does not help or
hurts: irrelevant retrieved context measurably degrades LLM
reasoning~\cite{distracted}, models underuse information positioned
mid-context~\cite{lostinmiddle}, and retrieval helps mainly where the
model's parametric knowledge is weak~\cite{whennottotrust} --- a pattern our
headroom results echo in the repair setting
(Section~\ref{sec:results-mechanism}). Our planted-corpus design contributes
a repair-specific mechanism: it separates \emph{whether the corpus contains a
sufficiently relevant example} (coverage) from \emph{whether the retriever
surfaces it} (ranking), which prior repair evaluations conflate.
